\clearpage
\phantomsection% 
\addcontentsline{toc}{chapter}{RINGKASAN}
\thispagestyle{fancy}

\begin{center}
	\large \bfseries \MakeUppercase{Ringkasan}\\
	\normalsize \normalfont {\thetitle}\\
	\normalsize \normalfont {\theauthor}\\
	\bigskip
	
	\normalsize \normalfont \justifying
	Perkembangan kecerdasan buatan generatif telah meningkatkan kemampuan pembuatan media sintetis, termasuk \textit{deepfake}, yang mampu menghasilkan citra atau video wajah dengan tingkat realisme tinggi. Kondisi ini menimbulkan tantangan dalam bidang keamanan informasi dan kepercayaan digital, karena konten \textit{deepfake} dapat digunakan untuk penyebaran informasi palsu, penyalahgunaan identitas, dan manipulasi opini publik. Deteksi \textit{deepfake} pada video menjadi semakin menantang karena data yang beredar di dunia nyata memiliki variasi kualitas visual, resolusi, pose, pencahayaan, dan karakteristik temporal yang tidak seragam. Oleh karena itu, diperlukan analisis terhadap pendekatan deteksi yang mampu bekerja pada data video \textit{deepfake} dalam kondisi \textit{in-the-wild}. \par

	Penelitian ini berfokus pada analisis komparatif antara pendekatan spasial dan pendekatan spasiotemporal untuk klasifikasi video \textit{deepfake}. Pendekatan spasial menganalisis informasi visual pada tingkat \textit{frame} tunggal, sedangkan pendekatan spasiotemporal memanfaatkan hubungan antar-\textit{frame} dalam satu klip video. Permasalahan utama dalam penelitian ini adalah belum jelasnya kontribusi relatif informasi spasial dan spasiotemporal dalam klasifikasi video \textit{deepfake} pada data \textit{in-the-wild}. Selain itu, efektivitas kedua pendekatan juga perlu dianalisis ketika karakteristik masukan dan konfigurasi pelatihan berubah, serta ketika model diterapkan pada \textit{dataset} eksternal dengan distribusi data yang berbeda. \par

	Rumusan masalah dalam penelitian ini mencakup bagaimana kontribusi pemanfaatan informasi spasiotemporal melalui VideoMAE dibandingkan informasi spasial melalui Vision Transformer dalam klasifikasi video \textit{deepfake} pada \textit{dataset} \textit{WildDeepfake}, bagaimana sensitivitas kedua pendekatan terhadap perubahan karakteristik masukan dan konfigurasi pelatihan, serta apakah efektivitas kedua pendekatan tetap konsisten ketika diterapkan pada \textit{dataset} eksternal \textit{Celeb-DF v2} dalam skenario \textit{zero-shot}. Berdasarkan rumusan masalah tersebut, penelitian ini bertujuan untuk menganalisis kontribusi informasi spasiotemporal dibandingkan informasi spasial, menganalisis sensitivitas pendekatan spasial dan spasiotemporal terhadap perubahan konfigurasi, serta mengevaluasi konsistensi efektivitas kedua pendekatan pada evaluasi lintas-\textit{dataset}. \par

	Metodologi penelitian dilakukan melalui beberapa tahapan, yaitu identifikasi masalah, studi literatur, pengumpulan data, \textit{preprocessing} data, perancangan model, pelatihan model, evaluasi model, serta analisis hasil. \textit{Dataset} utama yang digunakan adalah \textit{WildDeepfake}, sedangkan \textit{Celeb-DF v2} digunakan sebagai \textit{dataset} eksternal untuk evaluasi lintas-\textit{dataset} secara \textit{zero-shot}. Pada tahap \textit{preprocessing}, video wajah diproses menjadi klip berdurasi tetap sepanjang 46 \textit{frame}. Selanjutnya, dilakukan \textit{strided temporal sampling} untuk menghasilkan \textit{frame} masukan model. Data kemudian dibagi menjadi subset \textit{train}, validasi, dan \textit{test} dengan memperhatikan pencegahan kebocoran data antar-\textit{split}. Model spasial memproses setiap \textit{frame} dalam klip menggunakan Vision Transformer, kemudian hasil prediksi \textit{frame} diagregasi menjadi prediksi tingkat klip. Sementara itu, model spasiotemporal memproses seluruh \textit{frame} dalam satu klip secara langsung menggunakan VideoMAE. \par

	Hasil eksperimen \textit{baseline} menunjukkan bahwa model spasiotemporal memperoleh performa lebih tinggi dibandingkan model spasial pada \textit{dataset} \textit{WildDeepfake}. Model spasiotemporal BASE-ST memperoleh \textit{accuracy} sebesar 85,8\%, \textit{F1-score} sebesar 84,5\%, dan \textit{AUC} sebesar 92,3\%. Sementara itu, model spasial BASE-SP memperoleh \textit{accuracy} sebesar 79,2\%, \textit{F1-score} sebesar 76,7\%, dan \textit{AUC} sebesar 86,8\%. Hasil ini menunjukkan bahwa pemanfaatan hubungan antar-\textit{frame} pada VideoMAE memberikan kontribusi positif terhadap peningkatan performa klasifikasi video \textit{deepfake} pada evaluasi internal \textit{WildDeepfake}. Dengan demikian, informasi spasiotemporal mampu menyediakan evidensi tambahan yang tidak dimodelkan secara eksplisit oleh pendekatan spasial berbasis \textit{frame} tunggal. \par

	Pengujian konfigurasi pelatihan menunjukkan bahwa perubahan \textit{hyperparameter} memberikan pengaruh yang berbeda pada kedua pendekatan. \textit{Learning rate} yang terlalu besar dan augmentasi yang terlalu berat cenderung menurunkan performa model. \textit{Focal Loss} memberikan peningkatan pada model spasial, tetapi tidak memberikan peningkatan pada model spasiotemporal. \textit{Dropout} memberikan perbaikan kecil pada model spasial, tetapi menurunkan performa model spasiotemporal. Pada pengujian \textit{scheduler}, \textit{CosineAnnealingLR} dipilih sebagai konfigurasi lanjutan karena memberikan rata-rata \textit{accuracy} yang paling baik untuk kedua pendekatan. Selanjutnya, pengujian konfigurasi input menunjukkan bahwa resolusi, jumlah \textit{frame}, dan \textit{stride} memengaruhi performa model. Penurunan resolusi membatasi detail visual yang dapat dipelajari model, sedangkan perubahan jumlah \textit{frame} dan \textit{stride} memengaruhi cakupan serta kepadatan informasi temporal yang tersedia. \par

	Evaluasi lintas-\textit{dataset} pada \textit{Celeb-DF v2} menunjukkan pola yang berbeda dari evaluasi internal pada \textit{WildDeepfake}. Pada skenario \textit{zero-shot}, model spasial E6-SP memperoleh \textit{accuracy} sebesar 72,9\% dan \textit{AUC} sebesar 77,7\%, sedangkan model spasiotemporal E6-ST memperoleh \textit{accuracy} sebesar 64,2\% dan \textit{AUC} sebesar 72,6\%. Hasil ini menunjukkan bahwa efektivitas pendekatan spasiotemporal yang lebih tinggi pada \textit{dataset} \textit{WildDeepfake} tidak sepenuhnya konsisten ketika diterapkan pada \textit{dataset} eksternal. Dengan demikian, fitur temporal yang efektif pada \textit{dataset} sumber belum tentu memiliki tingkat transferabilitas yang sama ketika diuji pada distribusi data yang berbeda. \par

	Berdasarkan hasil penelitian, dapat disimpulkan bahwa pendekatan spasiotemporal berbasis VideoMAE memberikan kontribusi lebih besar pada evaluasi internal \textit{WildDeepfake} karena mampu memanfaatkan hubungan antar-\textit{frame} secara langsung. Namun, pendekatan spasial berbasis Vision Transformer menunjukkan konsistensi generalisasi yang lebih baik pada evaluasi lintas-\textit{dataset} terhadap \textit{Celeb-DF v2}. Oleh karena itu, efektivitas pendekatan spasial dan spasiotemporal perlu dilihat berdasarkan konteks evaluasinya. Pendekatan spasiotemporal lebih unggul pada evaluasi internal \textit{dataset} sumber, sedangkan pendekatan spasial lebih stabil pada skenario \textit{zero-shot} lintas-\textit{dataset}. \par

	\vfill
	
\end{center}
\clearpage
